\section{Security analysis}\label{sec:properties}

Each property of \S\ref{sec:threat} is a theorem of a Lean~4
model~\cite{lean4}. The model contains the guardian's \texttt{\_\_check\_auth}
as the function $\mathsf{check}$, its twelve checks in order with their error
codes; the circuit's relation as $\mathsf{Relation}(x, w)$, one field per
obligation C1 to C6 for public inputs $x = (\Cm, \digest_{hi}, \digest_{lo},
\mathit{root})$ and witness $w$; and Perch's controller for one account as a
state machine whose transitions cite Perch's transitions T1 to T9, with an
event log that records the configuration each event happened under.
$\mathsf{check}$ succeeds exactly when the facts of its twelve checks
hold\claim{check-iff-accepted}. SHA-256 and Poseidon2 are abstract
functions, and their breaks (a collision, or a Poseidon2 input that maps to
the empty leaf $0$) are named propositions. The header-level functions of
the circuit (which bytes are the \texttt{From} address, the Subject value,
the \texttt{d=} value; lowercasing; hex; the PKCS\#1 encoding) are abstract
too: computing them as \S\ref{sec:design-circuit} says is the circuit's
obligation, tested in \S\ref{sec:hardening} but not proven. No assumption is
an axiom: the build pins every theorem's axiom dependencies to Lean's
built-in three\claim{lean-axioms}.

\subsection{Soundness}

\begin{theorem}[Soundness of the email factor]\label{thm:soundness}
Let guardian $G$ have configuration $c$, and suppose
$\mathsf{check}(c, \mathit{env}, (S, R, \pi), \mathit{ctx})$ succeeds at
ledger $L$. If the configured verifier has a knowledge extractor (A4), there
are an entry point $f$, a digest $\digest$ and a witness $w$ such that:
\begin{enumerate}
  \item $\mathit{ctx} = [\,c.\mathit{controller}.f((\digest,))\,]$,
    $\Sha(S) = \digest$, and $S$ names $c$'s account, controller and
    network, an allowed action whose entry point is $f$, and a
    \texttt{valid\_until\_ledger} $\ge L$;
  \item $R$ is the pin's root with $L$ before its expiry, or a root the
    registry accepts at $L$;
  \item $w.H$ is well formed and carries a valid RSA-2048 PKCS\#1 v1.5
    signature $w.s$ under $w.n$ (C1, C2);
  \item $\leaff(w.\mathit{dom}, \kh(w.n))$ has a Merkle path to $R$ (C6);
  \item the \texttt{From} address $a$ of $w.H$ has at most 124 bytes and
    $\commitf(a, w.\code) = c.\Cm$ (C3);
  \item $\mathrm{lower}(\mathrm{domain}(a)) = \mathrm{lower}(\texttt{d=}) =
    w.\mathit{dom}$ (C4);
  \item the Subject of $w.H$ is the recovery phrase followed by
    $\mathrm{hex}(\blindf(\digest, w.\blind))$ (C5).
\end{enumerate}
If, moreover, every root the registry accepts is the root of a depth-16 tree
whose non-empty leaves are trusted (A5), and $c.\Cm$ commits to the enrolled
address $a^\ast$ while the pin holds only trusted leaves (A7), then
$a = a^\ast$ and the key's leaf is trusted, or Poseidon2 has a collision, or
Poseidon2 maps some input to $0$.
\end{theorem}
The first part is \textsf{EmailRecovery.email\_of\_check} and the whole is
\textsf{EmailRecovery.soundness}\claim{email-factor-soundness}. Through $S$,
the email binds the four items of \S\ref{sec:design-statement}. The step to
G1 is A1: a valid RSA-2048 signature under a trusted key of
$\mathrm{domain}(a^\ast)$ over a header whose \texttt{From} is $a^\ast$ was
produced by that domain's provider, which signed it for the holder of
$a^\ast$. That step is computational and outside the model. The Poseidon2
breaks matter because the key trees' empty leaf is $0$: a key leaf that
hashed to $0$ would stand in for an empty slot of any tree, which collision
resistance alone does not exclude, so A3 includes it. The proof uses two
lemmas about the spec's formulas: the commitment binds the address and the
code, and a Merkle path binds its leaf, each unless Poseidon2 collides.

\subsection{Privacy}

\begin{theorem}[Privacy as noninterference]\label{thm:privacy}
Let a guardian be enrolled from a secret $(a, \code, K)$, an address, an
account code and pinned keys. There is a function $f$ such that, for every
initial controller state $s_0$ and every sequence of transactions, the
chain's view of the run (the guardian's configuration, every verdict it
gives, and the controller's state and event log) equals
\[
  f\bigl(\commitf(a, \code),\ \mathrm{pinRoot}(\code, K),\ s_0,\ \text{transactions}\bigr).
\]
\end{theorem}
This is \textsf{EmailRecovery.privacy\_noninterference}, with the two-run form
\textsf{privacy\_two\_runs} and \textsf{publicInputs\_enroll}, which shows
the proof's public inputs are $\Cm$, the digest limbs and the
root\claim{address-privacy}. The theorem reduces privacy to the secrecy of
the preimages of $\Cm$ and the pin's root. G2, the computational statement
that a guardian's state and accepted approvals are indistinguishable between
two addresses at allowlisted providers when $\code$ is uniform and unknown,
follows only with A3's random-function model, A4's zero knowledge and A6,
and is not proven\claim{commitment-hiding}\claim{secret-generation}. What
remains visible is that the account has an email guardian, its allowed
actions, when each approval happened, and whether it used the pin or a
registry root. Off chain the provider sees $T$, which hides $\digest$ under
the same assumptions\claim{blinded-subject-hiding}.

\subsection{Replay and front-running resistance}

\begin{theorem}[Replay resistance]\label{thm:replay}
(a) If $\mathsf{check}$ accepts $(S, R, \pi)$ in contexts $\mathit{ctx}$, then
$\mathit{ctx} = [\,c.\mathit{controller}.f((\Sha(S),))\,]$; it accepts it in
exactly one context and for exactly one account, controller and network.
(b) Two statements with equal digests have the same account, network,
controller, action, attempt or request id and epoch, or SHA-256 collides.
(c) If the composed controller and guardian record the same approval in two
successful \texttt{submit\_guardian} calls, both are for one attempt, one
evidence domain and one epoch, or SHA-256 collides; a vote already recorded
is refused again. (d) The verdict ignores the authorization entry the
approval is wrapped in, and the controller's transitions take no submitter
input.
\end{theorem}
In Lean, parts (a) to (d) are \textsf{accepted\_names\_digest},
\textsf{accepted\_context\_unique} and
\textsf{accepted\_account\_unique};
\textsf{digest\_fixes}\claim{digest-fixes-statement};
\textsf{Perch.email\_vote\_replay}\claim{replay-resistance} with
\textsf{votes\_counted\_once} and \textsf{counted\_vote\_refused\_again}; and
\textsf{check\_payload\_irrelevant}\claim{guardian-payload-irrelevant} with
\textsf{emailAuth\_payload\_irrelevant}. A copy of an approval submitted by
anyone therefore records the vote the original would have recorded, for the
statement the original named. Part (b) uses two facts about Perch's encoding
as hypotheses: the encoding is injective, and the guardian's parser reads
back every field. They are tested, not proven: Perch's encoder reproduces
independent vectors, and the guardian parses every statement shape it
produces\claim{model-code-agreement}.

\subsection{One-factor safety}

\begin{theorem}[One-factor safety]\label{thm:onefactor}
In every state reachable from an unenrolled account by transactions whose only
recovery evidence is approvals by $G$ (owner-authorized transactions
allowed), every completed recovery of attempt $i$ is preceded in the log by
the authorization of $i$ at a ledger $L_a$ under a configuration $c$ with
mode \texttt{guardian-only}, quorum $1$ and $G$ among its guardians, and
$0 < c.\mathit{delay}$ and $L_a + c.\mathit{delay} \le L_\mathit{complete}$.
\end{theorem}
This is \textsf{EmailRecovery.Perch.one\_factor\_safety}\claim{one-factor-safety}.
It combines a statement about every trace, whatever its evidence (every
completion waits out a positive delay after its
authorization\claim{completion-waits-for-delay}), with one about traces whose
only evidence is $G$'s (every authorization is at quorum one, guardian-only);
so under a quorum of at least two, or in \texttt{combined} or
\texttt{zk-only} mode, email approvals alone authorize
nothing\claim{email-only-never-authorizes}. The same holds for the recovery
condition of a \texttt{Protected} reconfiguration or
upgrade\claim{email-alone-change-condition}, and during the delay a
\texttt{Loss} owner's cancellation of a live attempt
succeeds\claim{owner-can-cancel}. Under the wallet policy of
\S\ref{sec:design-policy}, stated as a predicate on recovery members, the
email guardian alone reaches the quorum only under \texttt{Loss}, with a
delay of at least 120{,}960 ledgers and a second alert channel
(\textsf{one\_factor\_under\_wallet\_policy})\claim{wallet-policy-one-factor}.
The chain accepts any configuration its owner applies, so this last part
binds Nido's wallet, not other wallets.

\subsection{Obligations}

Table~\ref{tab:obligations} lists, per property, what each component must
do and how it is established. The circuit's obligations and the contracts'
checks are tested; the composition is proven over the model; the
cryptographic steps are assumed.

\begin{table*}[t]
  \caption{Obligations of each property and their status.}
  \label{tab:obligations}
  \small
  \begin{tabularx}{\textwidth}{@{}l>{\raggedright\arraybackslash}p{0.12\textwidth}>{\raggedright\arraybackslash}X>{\raggedright\arraybackslash}p{0.17\textwidth}l>{\raggedright\arraybackslash}p{0.25\textwidth}@{}}
    \toprule
    Property & Circuit & Guardian and registry & Perch & Assumptions & Status \\
    \midrule
    1 Soundness & C1--C6 & checks 1--12; construction's circuit-id check; registry activation, grace, revocation and scope, with C6's hashing & statement built from controller state & A1, A2, A3, A4, A5, A7 & composition proven; C1--C6 and checks tested; registry rules tested; hashing cross-checked \\
    2 Privacy & public inputs exactly $(\Cm, \digest_{hi}, \digest_{lo}, \mathit{root})$ & stores $\Cm$ and a salted pin root; emits nothing & --- & A3, A4, A6 & noninterference proven; indistinguishability assumed \\
    3 Replay & C5 binds $\digest$ & checks 1, 2, 5, 6, 8, 12 & one vote per guardian per attempt and domain or per change digest; ids never reused; epochs grow & A2, A3, A4 & proven, with the encoding facts tested \\
    4 One factor & --- & the guardian is one address & quorum, delay, cancellation & --- & proven over all traces; wallet policy proven as a predicate and tested in code \\
    \bottomrule
  \end{tabularx}
\end{table*}

\subsection{Hardening against public issue classes}\label{sec:hardening}
Public audits of zkEmail's Circom and Noir
libraries~\cite{audit-zksecurity-2024,audit-yacademy,audit-zellic-2024,audit-matterlabs-2024,audit-veridise-2024,audit-consensys-2024}
and public issues in its repositories report recurring classes of
soundness bugs, mostly in header parsing rather than in DKIM itself; such
under-constrained circuits are a common class of SNARK
vulnerability~\cite{chaliasos-sok-snark}. The rules of
\S\ref{sec:design-circuit} refuse each publicly documented class;
Table~\ref{tab:hardening} in Appendix~\ref{app:hardening} lists the ten
classes, their public sources and the refusals. Each class has a negative test that names its source;
the test signs its header with the test key, so the RSA check passes and the
refusal comes from the check under test, and pins the exact failing
assertion\claim{circuit-issue-classes}; an attack on prover hints starts from
the honest hints and changes them as an attacker would. Three classes from the same reports do not
arise: partial-SHA state aliasing (no partial hashing is used), an unbound
Barrett parameter (it enters no constraint and is not hashed), and an
all-zero SHA output (a Circom template not used here).
Beyond the per-class tests, every adversarial email of the provider mocks,
given the witness most favourable to its claim, is refused at a named
assertion whenever its claim must not be provable\claim{circuit-refuses-attacks}
(Appendix~\ref{app:circuit}). Every unconstrained value is a hint that
constrained code re-checks, so a wrong hint can make an honest proof fail
but not make a false statement provable; the pinned compiler's public
advisories were reviewed against the constructs the circuit uses, and its
correctness beyond that review is assumed\claim{toolchain-correctness}.

\subsection{From the model to the code}\label{sec:gap}
The theorems are about the model; the model's agreement with the code is
tested, not proven. The model's $\mathsf{check}$ mirrors the guardian's
Rust \texttt{check} check for check, and each refusal is pinned by a unit
test\claim{guardian-checks}. The controller model mirrors Perch's
\texttt{contract.rs} transition by transition, and the guardian on Perch's
deployed WebAssembly is pinned by integration tests: one approval is one
vote, a resubmission is refused, approvals do not transfer to another attempt
or account, and email approvals alone never complete a
recovery\claim{guardian-on-perch}. The model's abstractions only add
behaviours, so a safety property of the model holds of the contracts:
ledger arithmetic is unbounded, storage never expires, and the target
derivation, the ZK adapter's verdict and the host's acceptance of an
authorization entry are inputs of a transaction. Two identifications remain:
configurations are compared by value where Perch compares
\texttt{config\_hash} (equal unless SHA-256 collides), and the two halves of
a completion are one transition. Not modelled are the account's own state
beyond its applied document, Perch's ZK membership pool, the circuit's
header parsing, and the semantics of Noir and Rust. The wallet-policy
predicate mirrors the wallet's builder rule for rule, and the builder is
tested at each rule's boundary. The registry's tree shape is tested and its
hashing cross-checked against an independent Noir
implementation\claim{registry-hash-parity}.
